\documentclass{article}
\usepackage{amsmath, amssymb, amsthm}

\setlength{\parindent}{0pt}

\newtheorem*{claim}{Claim}

\begin{document}

\begin{claim}
    There does not exist a positive integer $m$ such that
    \[
    (m+1)^3 + (m+2)^3 + \dots + (2m)^3
    \]
    is a perfect square.
\end{claim}

\begin{proof}
    Suppose for contradiction that there exists a positive integer $m$ such that
    \[
    (m+1)^3 + (m+2)^3 + \dots + (2m)^3 = n^2
    \]
    for some non-negative integer $n$.

    \medskip

    Using the identity $\sum_{i=1}^{N} i^3 = \left(\frac{N(N+1)}{2}\right)^2$, we get
    \[
    \begin{aligned}
        \sum_{i=m+1}^{2m} i^3
        &= \sum_{i=1}^{2m} i^3 - \sum_{i=1}^{m} i^3 \\
        &= \left(\frac{2m(2m+1)}{2}\right)^2 - \left(\frac{m(m+1)}{2}\right)^2 \\
        &= \left(\frac{4m^2+2m}{2}\right)^2 - \left(\frac{m(m+1)}{2}\right)^2 \\
        &= (2m^2+m)^2 - \left(\frac{m(m+1)}{2}\right)^2 \\
        &= \left(2m^2 + m + \frac{m(m+1)}{2}\right)\left(2m^2 + m - \frac{m(m+1)}{2}\right) \\
        &= \left(2m^2 + m + \frac{m^2}{2} + \frac{m}{2}\right)\left(2m^2 + m - \frac{m^2}{2} - \frac{m}{2}\right) \\
        &= \left(\frac{5}{2}m^2 + \frac{3}{2}m\right)\left(\frac{3}{2}m^2 + \frac{1}{2}m\right) \\
        &= \left(\frac{5m^2 + 3m}{2}\right)\left(\frac{3m^2 + m}{2}\right) \\
        &= \frac{m(5m+3)}{2} \cdot \frac{m(3m+1)}{2} \\
        &= \frac{m^2(5m+3)(3m+1)}{4}
    \end{aligned}
    \]
    Hence $m^2(5m+3)(3m+1) = 4n^2 = (2n)^2$.

    \medskip

    Since $x^2 \mid y^2$ implies $x \mid y$, we have $m \mid 2n$ for some positive integer $s$. Substituting gives 
    $m^2(5m+3)(3m+1) = m^2s^2$, hence
    \[
    (5m+3)(3m+1) = s^2
    \]

    \medskip

    We apply the Euclidean algorithm to find $\gcd(5m+3, 3m+1)$:
    \[
    \begin{aligned}
        \gcd(5m+3, 3m+1) 
        &= \gcd(3m+1, 5m+3 - 2(3m+1)) \\
        &= \gcd(3m+1, m-1) \\
        &= \gcd(m-1, 3m+1 - 3(m-1)) \\
        &= \gcd(m-1, 4)
    \end{aligned}
    \]
    so $d := \gcd(5m+3, 3m+1)$ divides $4$, i.e.\ $d \in \{1, 2, 4\}$. Write $5m+3 = d\alpha$ and $3m+1 = d\beta$ 
    where $\gcd(\alpha, \beta) = 1$. Then $(5m+3)(3m+1) = s^2$ gives $d^2\alpha\beta = s^2$, so $\alpha\beta$ is a 
    perfect square. Since $\gcd(\alpha,\beta)=1$, no prime divides both, so each prime appears in only one factor. 
    Hence each prime appears with an even exponent in each factor, so both $\alpha$ and $\beta$ are perfect squares. 
    We then case-split on $d \in \{1, 2, 4\}$.

    \medskip
    \textbf{Case 1: $d = 1$.} Then $5m+3 = a^2$ and $3m+1 = b^2$ for integers $a, b$. Eliminating $m$ gives $3a^2 - 
    5b^2 = 4$. Reducing mod $5$ gives $3a^2 \equiv 4 \pmod{5}$. However, $a^2 \equiv 0, 1, 4 \pmod{5}$, so $3a^2 
    \equiv 0, 3, 2 \pmod{5}$, and $4$ is not among these residues. A contradiction.

    \medskip
    \textbf{Case 2: $d = 2$.} Then $\frac{5m+3}{2} = a^2$ and $\frac{3m+1}{2} = b^2$ for integers $a, b$. Eliminating $m$ 
    gives $3a^2 - 5b^2 = 2$. Reducing mod $3$ gives $-5b^2 \equiv 2 \pmod{3}$. Since $-5 \equiv 1 
    \pmod{3}$, this simplifies to $b^2 \equiv 2 \pmod{3}$. However, $b^2 \equiv 0, 1 \pmod{3}$, and $2$ is not among these 
    residues. A contradiction.

    \medskip
    \textbf{Case 3: $d = 4$.} Then $\frac{5m+3}{4} = a^2$ and $\frac{3m+1}{4} = b^2$ for integers $a, b$. Eliminating $m$ 
    gives $3a^2 - 5b^2 = 1$. Reducing mod $5$ gives $3a^2 \equiv 1 \pmod{5}$. However, $a^2 \equiv 0, 1, 4 \pmod{5}$, so 
    $3a^2 \equiv 0, 3, 2 \pmod{5}$, and $1$ is not among these residues. A contradiction.

    \medskip

    In all three cases we reach a contradiction, so no such $m$ exists.
\end{proof}

\end{document}
